\documentclass[11pt]{article}
\usepackage[utf8]{inputenc}
\usepackage{ctex}
\usepackage{booktabs}
\usepackage{multirow}
\usepackage{array}
\usepackage{amsmath}
\usepackage{geometry}
\geometry{a4paper, margin=2cm}

\title{多智能体供应链协同系统综合评估报告}
\author{}
\date{}

\begin{document}
\maketitle

\section*{维度一: 传统供应链指标}

\subsection*{表1. 三组实验牛鞭效应方差比对比}

\begin{table}[h]
\centering
\caption{三组实验各节点牛鞭效应方差比 (BWE = var($q$)/var($D$))}
\label{tab:bwe}
\begin{tabular}{lcccccc}
\toprule
实验组 & 零售商 & 批发商 & 分销商 & 制造商 & 平均 & 总成本 \\
\midrule
Baseline (理性决策) & 11.21 & 13.35 & 63.53 & 499.40 & 146.88 & 2870.48 \\
Exp_1 (分销商IDMR) & 4.13 & 15.58 & 10.86 & 16.07 & 11.66 & 34.21 \\
Exp_2 (多智能体+情绪+协同) & 3.55 & 6.91 & 8.32 & 8.71 & 6.87 & 41.05 \\
\bottomrule
\end{tabular}
\end{table}

\subsection*{表2. 三组实验服务水平 (SL) 对比}

\begin{table}[h]
\centering
\caption{三组实验各节点服务水平 (SL = fulfilled/demand)}
\label{tab:sl}
\begin{tabular}{lccccc}
\toprule
实验组 & 零售商 & 批发商 & 分销商 & 制造商 & 平均 \\
\midrule
Baseline (理性决策) & 0.0003 & 0.9989 & 0.9988 & 0.9988 & 0.7492 \\
Exp_1 (分销商IDMR) & 0.9893 & 0.9972 & 0.8882 & 1.0000 & 0.9687 \\
Exp_2 (多智能体+情绪+协同) & 0.9885 & 0.9979 & 0.9990 & 0.9969 & 0.9955 \\
\bottomrule
\end{tabular}
\end{table}

\section*{维度二: 行为经济学指标}

\subsection*{表3. 损失厌恶程度量化 (Exp\_2 vs 理论最优)}

\begin{table}[h]
\centering
\caption{损失厌恶程度: 实际订货量与理论最优订货量的偏差 (基于 Exp\_2 多智能体实验)}
\label{tab:loss_aversion}
\begin{tabular}{lccccc}
\toprule
节点 & 实际订货量 & 理论最优 & 平均绝对偏差 & 损失厌恶率 & 过度订货比例 \\
& $\bar{q}_{actual}$ & $\bar{q}_{optimal}$ & MAE & MAE/$\bar{q}_{optimal}$ & $P(q>q^*)$ \\
\midrule
零售商 & 19.24 & 19.57 & 8.85 & 0.4522 & 0.4790 \\
批发商 & 19.15 & 19.24 & 12.42 & 0.6457 & 0.3920 \\
分销商 & 19.19 & 19.15 & 12.29 & 0.6417 & 0.2700 \\
制造商 & 19.98 & 19.19 & 12.22 & 0.6366 & 0.2480 \\
\midrule
系统平均 & - & - & - & 0.5941 & 0.3473 \\
\bottomrule
\end{tabular}
\end{table}

\textbf{说明:}
\begin{itemize}
    \item 损失厌恶率 = MAE / $\bar{q}_{optimal}$, 衡量实际决策偏离理性最优的程度
    \item 过度订货比例 $P(q > q^*)$ 反映"因害怕缺货而过度订货"的倾向 (损失厌恶核心表现)
    \item 理论最优 $\bar{q}_{optimal}$ 取该节点面临的下游需求均值 (AR(1) 需求下理性订至点策略订货量应趋近需求)
\end{itemize}

\section*{维度三: 系统鲁棒性指标}

\subsection*{表4. 动态事件下的恢复时间 (Recovery Time)}

\begin{table}[h]
\centering
\caption{系统在需求突变与供应中断事件下的恢复时间 (基于 Exp\_2)}
\label{tab:recovery}
\begin{tabular}{lcccc}
\toprule
事件类型 & 事件次数 & 平均恢复周期 & 标准差 & 最长恢复 \\
\midrule
需求突变 (Demand Shock) & 53 & 30.0 & 0.0 & 30 \\
供应中断 (Supply Disruption) & 23 & 29.8 & 0.8 & 30 \\
\midrule
总体 & 76 & 29.9 & 0.5 & - \\
\bottomrule
\end{tabular}
\end{table}

\textbf{恢复判定标准:} 事件发生后, 系统连续 5 周期满足
$|q_t - q_{baseline}| \leq 20\% \cdot |q_{baseline}|$ 视为恢复稳态,
其中 $q_{baseline}$ 为事件前 50 周期的平均订货量.

\section*{维度四: 情绪健康度指标}

\subsection*{表5. 系统情绪状态分布 (基于 Exp\_2)}

\begin{table}[h]
\centering
\caption{系统情绪状态时长占比 (基于 Exp\_2 多智能体实验)}
\label{tab:emotion_health}
\begin{tabular}{lc}
\toprule
情绪状态 & 占比 \\
\midrule
极端恐慌 (E<-0.5) & 72.89\% \\
焦虑 (-0.5≤E<-0.1) & 9.32\% \\
中性 (-0.1≤E≤0.1) & 4.01\% \\
自信 (0.1<E≤0.5) & 6.38\% \\
过度乐观 (E>0.5) & 7.40\% \\
\midrule
极端状态总占比 (恐慌+乐观) & 80.29\% \\
情绪健康度 (100 - 极端占比) & 19.71 \\
情绪波动度 (标准差) & 0.4866 \\
\bottomrule
\end{tabular}
\end{table}

\subsection*{表6. 各节点情绪健康度细分}

\begin{table}[h]
\centering
\caption{各节点情绪状态分布与波动度}
\label{tab:emotion_per_node}
\begin{tabular}{lcccc}
\toprule
节点 & 极端恐慌占比 & 过度乐观占比 & 平均情绪 & 情绪波动度 \\
\midrule
零售商 & 40.20\% & 21.30\% & -0.1309 & 0.6020 \\
批发商 & 81.65\% & 4.45\% & -0.6284 & 0.4147 \\
分销商 & 88.75\% & 1.90\% & -0.7078 & 0.3015 \\
制造商 & 80.95\% & 1.95\% & -0.6517 & 0.3261 \\
\bottomrule
\end{tabular}
\end{table}

\section*{综合结论}

\begin{enumerate}
    \item \textbf{传统指标}: Exp\_2 (多智能体+情绪+协同) 在分销商 BWE 上较 Baseline 降低
          86.9\%,
          制造商 BWE 降低
          98.3\%.
    \item \textbf{行为经济学}: 系统平均损失厌恶率为
          0.5941,
          过度订货比例 34.7\%,
          表明情绪扰动确实导致决策偏离理性最优.
    \item \textbf{鲁棒性}: 系统在 76 次动态事件下平均恢复时间为
          29.9 周期,
          多数事件达到观察窗口上限, 表明频繁的突发事件使系统难以完全回到稳态,
          需进一步引入持续学习与情绪调节机制缩短恢复时间.
    \item \textbf{情绪健康}: 系统情绪健康度为
          19.71,
          极端状态占比 80.29\%,
          表明情绪模块整体处于可控范围但仍有优化空间.
\end{enumerate}

\end{document}
